\documentclass[aps,prd,preprint,superscriptaddress,amsmath,amssymb]{revtex4-2}
\usepackage{graphicx}
\usepackage{bm}
\usepackage{dcolumn}
\usepackage{amsthm}
\usepackage{booktabs}
\usepackage{color}

\newtheorem{postulate}{Postulate}[section]
\newtheorem{theorem}{Theorem}[section]

\begin{document}

\title{Divergence, Sink Contraction, and the Half-Phase Geometry:\\
Charge, the Antigravity Field, the Photon, and Spin~$1/2$ from Flowing Space}

\author{Yuan-Jie Liu}
\email{yuanjieliu65@gmail.com}

\date{\today}

\begin{abstract}
We extend the flowing-space framework in two geometric directions and
join them into one chain.  (i)~\emph{The half-phase of the critical
sheet}: the inversion fixed point of the envelope reflection
$w\mapsto e/w$ sits at $|w|=\sqrt{e}$, so $\ln(\sqrt e)=1/2$; the square
root is the half-step operator of the multiplicative group, and in phase
space the same half-step is the half-angle phase
$\exp(i\theta/2)$, whose closure factor is $-1$ --- the explicit
fermionic witness of spin~$1/2$.  (ii)~\emph{Charge as flow
divergence}: positive charge is a source (positive divergence),
negative charge a sink (negative divergence), mutually exclusive,
flipped under $\rho\to-\rho$, neutral at zero.  A time-varying source
forces the space field to vary (CR1); a varying field activates the
nuclear channel $m\,dC$ of the four-force decomposition (CR2); and a
sink --- a converging flow --- is geometrically a flattening: the
inward contraction $v_i\mapsto(1-\lambda)v_i+\lambda\bar v$ reduces the
fluctuation energy $Q_A$ by the square factor $(1-\lambda)^2$ (CR9,
exact discrete theorem).  Joining the framework's own $\mu$-bridge
(one flattening gives $\mu=1$), the chain closes: accelerated charge
$\Rightarrow$ antigravity field $\Rightarrow$ the photon (mass zero and
neutral).  As a device-level application, the self-flattening of
circulating electrons shortens the approach to the working window
$\mu=0.999$ from $132$ to $85$ steps and reaches it with halved
external RMF drive.  All statements are formalized in Lean~4 with zero
\texttt{sorry}; the honest status is stated: $\lambda$ (the sink
contraction rate) and $\delta,\varepsilon$ (the sink energy scale) are
inputs, and the continuum 3D divergence-to-functional correspondence is
not yet formalized.
\end{abstract}

\maketitle

\begin{quote}
\itshape
``Flowing water does not fight; it yields, and in yielding it reaches
everywhere.''
\par
\hfill --- after \textit{Daodejing}
\end{quote}

\section{Introduction}
\label{sec:intro}

The flowing-space framework \cite{ProjectionPhysics,HibsPhysics}
adopts a single primitive --- \emph{flow} --- and derives mass,
information, causality, force, the photon, and the quantum relations
from a small set of postulates: vector light speed, three anchored
directions, and the counting of space-displacement strands.  The
present work adds two geometric structures that were absent there:
the \emph{half-phase geometry} of the critical sheet, and the
\emph{divergence structure} of charge as source and sink.

The two additions are independent in content but share a common
theorem-level fact: a quadratic quantity is reduced by a square
factor.  For spin, the phase closure factor of the half-angle map is
$-1=(-1)^2$-free: the half step of the variable $\theta$ is the
square root in logarithmic coordinates, and the fermionic witness is
$\exp(i\theta/2)$.  For charge, the fluctuation energy of a
converging flow is reduced by $(1-\lambda)^2$ per contraction step;
the sink is a self-flattening structure, and a sink at full strength
($\lambda\to 1$) brings the fluctuation to zero, mass to zero, and
charge to neutrality --- the photon.

\section{Postulates of the flowing-space framework}
\label{sec:postulates}

We recall the postulates needed here; full derivations are in
Ref.~\cite{ProjectionPhysics}.

\begin{postulate}[Vector light speed]
\label{post:sls}
Space itself flows at the vector speed $C$ with $|C|=c$; the
direction of $C$ is changeable.
\end{postulate}

\begin{postulate}[Three anchored directions]
\label{post:three}
Space has three independent orthogonal directions; a physical object
is an anchored pattern of the flow, with anchor strength
$s\ge 0$.  The anchored mass is
\begin{equation}
\label{eq:anchor}
m^2_{\rm eff}\;=\;s^2(1-\mu)^2 ,
\end{equation}
where $\mu\in[0,1]$ is the antigravity strength
(flattening of the flow gradient); $\mu=1$ gives $m_{\rm eff}=0$
(MassCancellation AMC1).
\end{postulate}

\begin{postulate}[Four forces as Leibniz channels]
\label{post:four}
The flow momentum is $P=m(C-v)$; its time derivative decomposes into
four channels (GQF2):
\begin{equation}
\label{eq:four}
\frac{d}{dt}\bigl[m(C-v)\bigr]
\;=\;\underbrace{\dot m\,C}_{\text{electric}}
+\underbrace{m\,\dot C}_{\text{nuclear}}
-\underbrace{\dot m\,v}_{\text{magnetic}}
-\underbrace{m\,\dot v}_{\text{gravity}} .
\end{equation}
\end{postulate}

\section{Charge as flow divergence}
\label{sec:charge}

Charge is the divergence of the space-displacement flow
(GQF4; VibrationChargeFlow CF1--CF5).

\begin{theorem}[Source and sink]
\label{def:srcsink}
For a local divergence scalar $\rho$:
\begin{align}
\text{PositiveSource}(\rho) &:\Longleftrightarrow 0<\rho ,\\
\text{NegativeSink}(\rho)   &:\Longleftrightarrow \rho<0 .
\end{align}
Positive charge is a source (diverging displacement), negative charge
a sink (converging displacement); both obey the right-handed helical
motion of the flow, the sign of the divergence distinguishing them.
\end{theorem}

\begin{theorem}[Source--sink exclusivity; CF1--CF4]
\label{thm:srcsink}
(i)~A source and a sink cannot hold simultaneously.  (ii)~Charge
conjugation is the sign flip: $0<\rho \iff \rho<0$ after
$\rho\mapsto-\rho$ (source and sink are exchanged).  (iii)~Zero
divergence is neutral: neither source nor sink.  (iv)~The magnitude is
preserved under the flip: $|{-}\rho|=|\rho|$.
\end{theorem}

The continuity side --- the coupled extraction of the charge $e$, the
divergence dynamics in 3D, and charge normalization --- remains open
(honestly registered in the repository).

\section{Accelerated charge implies the field varies; the nuclear
channel activates}
\label{sec:accel}

\begin{theorem}[CR1; time-varying source forces a varying field]
\label{thm:cr1}
In the sourced wave equation
$\partial_t^2 C=c^2\partial_x^2 C+J$ (MS5), a static field implies a
static source: if $\partial_t C\equiv 0$, then $J(t{+}1)=J(t)$ for all
$t$.  Hence a time-varying source (an accelerated charge) forbids a
static space field.
\end{theorem}

\begin{theorem}[CR2; the nuclear channel activates]
\label{thm:cr2}
With $m\neq 0$ and $\dot C\neq 0$, the nuclear channel of
Eq.~(\ref{eq:four}) is non-zero: $m\,\dot C\neq 0$.
\end{theorem}

The missing link between the nuclear channel and the antigravity
parameter $\mu$ is not invented here: the framework already contains
its own $\mu$-bridge.

\begin{theorem}[CR3; one flattening gives $\mu=1$]
\label{thm:cr3}
Let $Q_A(v)=\sum_{i\in A}(v_i-\bar v_A)^2$ be the fluctuation energy,
and $\eta=1-Q_A(v)/Q_A(v_0)$ the flattening progress (MuFieldCoupling
TD11).  The $\mu$ update is $\mu'=\mu+\eta(1-\mu)$
(PlasmaDynamics muStep).  Flattening once gives
$\mu'=\mu+1\cdot(1-\mu)=1$.
\end{theorem}

\begin{theorem}[CR4; $\mu=1$ is the photon endpoint]
\label{thm:cr4}
At $\mu=1$, the anchored mass vanishes
(Eq.~\ref{eq:anchor} with $1-\mu=0$) and the divergence is neutral
(CF3).  The object is massless and chargeless --- the photon.
\end{theorem}

\begin{theorem}[CR5; closure of the chain]
\label{thm:cr5}
One flattening applied to any region $A$ (any anchored substance $s$)
yields a photon: mass zero and neutrality.
\end{theorem}

\section{Convergence is flattening: sink contraction}
\label{sec:sink}

Why is the \emph{negative} charge the natural antigravity carrier?
The geometric answer: a sink is a converging flow, and convergence is
flattening.

\begin{theorem}[Sink contraction; CR9]
\label{def:sinkcontract}
For a region $A$, a field $v:A\to\mathbb R$, and $\lambda\in[0,1]$,
\[
\text{sinkContract}(A,v,\lambda)(i)
\;=\;\begin{cases}
(1-\lambda)\,v_i+\lambda\,\bar v_A & i\in A,\\
v_i & i\notin A .
\end{cases}
\]
Each point of $A$ is pulled toward the region mean by a fraction
$\lambda$ --- the discrete image of ``the negative charge absorbs the
surrounding flow toward itself.''
\end{theorem}

\begin{theorem}[CR9; square decay of fluctuation]
\label{thm:cr9}
The contraction preserves the region mean and scales the fluctuation
energy by $(1-\lambda)^2$:
\begin{equation}
\label{eq:cr9}
Q_A\bigl(\text{sinkContract}(A,v,\lambda)\bigr)
\;=\;(1-\lambda)^2\,Q_A(v).
\end{equation}
For $0<\lambda<1$ the decay is monotone and exact; $\lambda=1$ is a
single flattening (fluctuation to zero).  The sink is a
self-flattening operator family.
\end{theorem}

\begin{theorem}[CR10; strict attraction pointwise]
\label{thm:cr10}
For $0<\lambda<1$ and $v_i\neq\bar v_A$,
\[
\bigl|\text{sinkContract}(A,v,\lambda)(i)-\bar v_A\bigr|
\;<\;\bigl|v_i-\bar v_A\bigr| ,
\]
each point is strictly attracted to the mean --- ``the sink pulls the
surrounding flow to itself,'' pointwise.
\end{theorem}

The energy side of the sink is already present in the framework:
information sinking into the absorbing branch $iR$ releases
$\varphi(R)-\varphi(iR)=\delta+\varepsilon>0$ (HiddenQFT HQ9a), so a
converging flow carries its own flattening power as a candidate
(CR8).  This is a semantic correspondence (interface), not a theorem
of the divergence dynamics.

\section{The half-phase geometry: spin~$1/2$ from the critical sheet}
\label{sec:halfphase}

The second geometric addition concerns the value $1/2$ itself.

\subsection{The inversion fixed point}

The Riemann--Hilbert envelope of the framework
\cite{RiemannHIBS} carries the reflection $s\mapsto 1-s$; in the
logarithmic coordinate $w=e^s$ this is the inversion $w\mapsto e/w$.
Its fixed set is the circle $|w|=\sqrt e$ (envelope inversion fixed
circle).  Taking logarithms:
\begin{equation}
\label{eq:half}
\ln(\sqrt e)\;=\;\tfrac12\ln e\;=\;\tfrac12 .
\end{equation}

\begin{theorem}[CS2; inversion fixed iff log half]
\label{thm:cs2}
For $r>0$,
\[
r=e/r \;\Longleftrightarrow\; \ln r=\tfrac12 .
\]
The fixed point of the reflection must sit at $1/2$ in the
logarithmic coordinate --- because $\ln e=1$ and the fixed point of
$s\mapsto 1-s$ is $1/2$.
\end{theorem}

\subsection{The half-angle phase is fermionic}

The square root is the half-step operator of the multiplicative group;
in the phase space of a vibration, the same half-step is the half-angle
phase.

\begin{theorem}[CS4; half-angle phase has closure factor $-1$]
\label{thm:cs4}
$\theta\mapsto\exp(i\theta/2)$ satisfies
\[
f(\theta+2\pi)\;=\;-1\cdot f(\theta),
\]
i.e.\ it is a fermionic closure witness (VibrationStatistics
ClosureFactor with $\sigma=-1$): $2\pi$ changes sign, $4\pi$
restores.  The half-angle phase is the explicit witness of
spin~$1/2$ in the phase space of a vibration.
\end{theorem}

\subsection{Energy and mass}

With $E\propto m^2$ (mass--energy quadrature), logarithms are
linear: $\ln E=2\ln m$.  At the critical sheet, $\ln(\sqrt e)=1/2$
means $m=\sqrt e$, $E=e$; the factor $2$ in $\ln E=2\ln m$ is the
inverse of the spin half-step --- energy takes two half-steps to
complete one cycle of $m$.

\begin{theorem}[CS5; logarithmic $2{:}1$]
\label{thm:cs5}
$\ln e=2\ln(\sqrt e)$.
\end{theorem}

\section{Device application: two-flow-ring optimization}
\label{sec:device}

The sink-flattening theorem (CR9) enters the $\mu$ dynamics of the
two-flow-ring device \cite{DesignTwoFlowRing} as a self-flattening
gain:
\begin{equation}
\label{eq:etasink}
\eta_{\text{sink}}(\lambda)=2\lambda-\lambda^2 ,
\end{equation}
the fraction of fluctuation removed by one contraction, added to the
external RMF gain and decaying with the residual
$\eta_{\text{sink}}\cdot(1-\mu)$ (the flattening loses its object as
the region flattens; TD18).  The comparison
(evaluated numerically; honest candidate, $\lambda$ an input):

\begin{table}[h]
\caption{$\mu$ trajectory to the working window $\mu_{\rm work}=0.999$
(FC11 ceiling $0.999780568$).}
\label{tab:sinkmu}
\begin{ruledtabular}
\begin{tabular}{lcc}
& Pure external RMF & External $+$ self-flattening ($\lambda=0.3$) \\
\hline
Steps to $\mu=0.999$ (CR1 criterion) & $132$ & $85$ \\
Steps with halved external drive & $267$ & $150$ \\
\end{tabular}
\end{ruledtabular}
\end{table}

The self-flattening shortens the approach to the working window by
$\sim35\%$ and reaches it with halved external RMF drive --- the
sink's own converging geometry contributes to the $\mu$ production
(supply-side, CR8 candidate).  The negative result is registered
honestly: a \emph{constant} $\eta_{\text{sink}}$ violates the FC11
ceiling (pushes $\mu$ to $1$); the decaying form is required.  With a
constant external drive, $\mu$ still tends to $1$ (a known TM-family
behavior); stopping inside the window requires feedback shutoff
(TD19), a separate matter.

\section{Numerical verification}
\label{sec:numerics}

All theorems of the chain are verified in Lean~4 with zero
\texttt{sorry} (4348 jobs, full build green) and numerically at
machine precision where a numeric witness exists:

\begin{itemize}
\item 3D $21^3$ lattice, negative-divergence field $v=-r$: every
interior point has $\operatorname{div} v<0$; initial fluctuation
$Q=10187$; 200 sink-contraction steps drive $Q\to 3.5\times10^{-61}$
with $|{\rm div}|\to 0$, monotonically (verify\_sink\_contract\_3d).
\item Sink contraction scaling $Q\to(1-\lambda)^2Q$ reproduced exactly
(CR9 witness); strict pointwise attraction (CR10) monotone in
$\lambda$.
\item Two-flow-ring comparison (Tab.~\ref{tab:sinkmu}): $132\to85$
steps at equal external drive; $267\to150$ with halved external
drive; $\eta_{\text{sink}}(\lambda)$ strictly increasing in $\lambda$
(verify\_sink\_mu\_optimization).
\end{itemize}

\section{Honest assessment}
\label{sec:honest}

\paragraph{Claimed.}
(i)~Charge is the divergence of the flow: source $=$ positive
divergence, sink $=$ negative divergence, with the four algebraic
facts (exclusivity, flip, neutrality, magnitude) formalized.
(ii)~The chain ``accelerated charge $\Rightarrow$ varying field
$\Rightarrow$ nuclear channel $\Rightarrow$ one flattening
$\Rightarrow$ $\mu=1$ $\Rightarrow$ photon'' is a closed conditional
chain of theorems: \emph{if} a flattening is applied, any anchored
substance becomes a photon.  (iii)~A sink is geometrically a
self-flattening operator: the exact discrete theorem
$Q_A\to(1-\lambda)^2 Q_A$ under sink contraction.  (iv)~The half-step
geometry locates spin~$1/2$: $\ln(\sqrt e)=1/2$ is the fixed point of
the envelope reflection in the logarithmic coordinate, and the
half-angle phase is the explicit fermionic closure witness.
(v)~All statements are formalized in Lean~4 with zero
\texttt{sorry}; the device comparison is numerically verified.

\paragraph{Not claimed.}
(i)~The bridge ``nuclear channel $\Rightarrow$ flattening action''
is \emph{not} derived: who applies the flattening (the physical
origin of the flattening power) is the second-input gap, unchanged.
(ii)~$\lambda$ (the sink contraction rate) and $\delta,\varepsilon$
(the sink energy scale) are inputs, not derived values; 
$\eta_{\text{sink}}=2\lambda-\lambda^2$ is a candidate implementation,
not the unique mapping.  (iii)~The continuum 3D correspondence
``negative divergence $\Rightarrow$ fluctuation functional decreases''
is not formalized (the discrete lattice theorems are the algebraic
seed).  (iv)~Spin~$1/2$ is \emph{located} geometrically, not
\emph{derived} as a value: the other half of the
spin--statistics theorem (phase $\Rightarrow$ operator anticommutation)
requires field-operator language absent from the framework --- the
result is half a theorem.  (v)~No new falsifiable prediction beyond
the standard framework; all numerical values reproduce known physics.

\section{Conclusion}
\label{sec:conclusion}

Two geometric structures extend the flowing-space framework.  The
divergence structure of charge makes the sink a self-flattening
operator, closing --- conditionally --- the chain from accelerated
charge to the photon and giving the two-flow-ring device a
self-flattening gain that shortens the approach to the working window
by a third.  The half-phase geometry locates spin~$1/2$ at the
envelope-reflection fixed point $\ln(\sqrt e)=1/2$, with the
half-angle phase as its explicit fermionic witness.  Both additions
are theorem-exact where they are exact, and honestly open where the
inputs remain: $\lambda$, $\delta$, $\varepsilon$, and the continuum
3D correspondence.

\begin{thebibliography}{9}

\bibitem{ProjectionPhysics}
Y.-J.~Liu \emph{et al.}, ``The geometric description of physics at the
cosmic scale,'' companion manuscript (paper/), repository
\texttt{logos-42/Hibs-Physics}.

\bibitem{HibsPhysics}
\texttt{logos-42/Hibs-Physics} repository: Lean~4 formalization and
numerical verification, \url{https://github.com/logos-42/Hibs-Physics}.

\bibitem{RiemannHIBS}
\texttt{logos-42/RiemannHIBS}: hidden-number envelope with the
reflection $s\mapsto 1-s$ and the inversion fixed circle $|w|=\sqrt e$.

\bibitem{DesignTwoFlowRing}
Y.-J.~Liu, ``Two-flow-ring design: gap-tolerant feasible region,''
\emph{Design manual} of the repository (docs/wiki/design-two-flow-ring.md),
and the companion device project \texttt{logos-42/HushFusion}.

\end{thebibliography}

\end{document}